% page 555 du fac-simile (image p-554 du scan)
% bloc 167.6 x 250.6 mm ; marge gauche 34.4 mm
\pgc{12.700mm}{0.000mm}{
\l{}
\l{}
\l{}
\l{}
\l{}
\l{\sou{spleno}. (\sou{anat.)}\cel{2}Vicero sponjatra, saturita ye sango, en la}
\l{\cel{3}hipokondro sinistra, di qua la rolo esas ne-bone savata -}
\l{\cel{3}e quan, olim, onu egardis kom la fonto di melankolio. - DEFRS}
\l{}
\l{\sou{splenalgio}. (\sou{patol}.) Doloro ye la spleno. - DEFIRS.}
\l{}
\l{\sou{splendida}. Di qua la luxozeso brilanta esas plena ye grandeso.}
\l{\cel{3}DEFIS.}
\l{}
\l{\sou{splenito}. (\sou{patol}.)Inflamuro di la spleno. - DEFI.}
\l{}
\l{\sou{splinar.(netrans.}) Afektesar da ta hipokondrio qua karakterizis}
\l{\cel{3}la Britaniani, e su manifestis per la enoyo pri omno. - DEFRS.}
\l{}
\l{\sou{splinto.}\cel{2}(\sou{tekn.)}\cel{2}Stifto qua asemblas la parti di charniri, di}
\l{\cel{3}bukli, e c. - DR.}
\l{}
\l{\sou{splis}ar. (\sou{trans.}) (\sou{tekn.}) Ligar (du kordi), extremajo an extre-}
\l{\cel{3}majo, interplektante la toroni. - DEF.}
\l{}
\l{\sou{splito.}\cel{2}I. Fragmento forlansita da korpo qua explozas. - II.}
\l{\cel{3}Mikra peceto de ligno, de metalo, e c., tenua ed akuta, qua}
\l{\cel{3}penetras la pelo acidente. - DE.}
\l{}
\l{\sou{spoko}. (\sou{tekn.})\cel{2}Singla de la parti di roto qui iras de la nabo}
\l{\cel{3}a la felgo. - E.}
\l{}
\l{\sou{spoliar}. (\sou{trans}.)Deprenar (de ulu), per ruzo o per violento,}
\l{\cel{3}por proprigar lu a su (to quo apartenis a lu). - DEFIS.}
\l{}
\l{\sou{spondeo}. (\sou{versifado}).Pedo longa ye du silabi. - DEFIRS.}
\l{}
\l{\sou{sponjo}. I. (\sou{zool.}) Zoofito, qua konsistas ek amaso de tisui fi-}
\l{\cel{3}broza, plu o min densa, flexebla, elastika, tre lejera,}
\l{\cel{3}imbibebla, e qua, kande vivanta, indutesas de materio gelatina-}
\l{\cel{3}tra mi-fluida, iritebla ed efemera. - L. spongia.- II. La}
\l{\cel{3}korpo di ta zoofito, uzata diversa-rolo, pro lua karakterizivo:}
\l{\cel{3}absorbar la liquidi, e lasar li ekirar lor preseso mem tre}
\l{\cel{3}poka. - EFIS.}
\l{}
\l{\sou{spontana}. I. Quan onu agas o facas ipse, propra-move, sen impuls-}
\l{\cel{3}eso da forco aplikata de extere. - II. Qua aspektas agata o}
\l{\cel{3}facata propra-move - di qua la kauzo ne esas perceptebla.}
\l{\cel{3}- DEFIS.}
\l{}
\l{\sou{sporo}. (bot.)\cel{2}Korpuskulo genitera (di ula planti kriptogama).-}
\l{\cel{3}DEFIRS.}
\l{}
\l{\sou{sporadika.}\cel{1}I. (geol.)\cel{2}Dispersita. - II. (\sou{patol.)}\cel{1}Qua afektas}
\l{\cel{3}individui quivivas ne-grupe - qua ne esas epidemia.- DEFIRS.}
\l{}
\l{\sou{sporangio.(bot.)}\cel{1}Quaza sako mikra, qua kontenas la spori, che}
\l{\cel{3}la kriptogami. - DEFIS.}
}
